%-*-tex-*-
\ifundefined{writestatus} \input status \relax \fi %
\chcode{fonts}
\def\cqu{}


%========  Special fonts for this chapter ========

\chapterhead{fonts}{FONTS}
\intex\ has available to it a large number of fonts.\footnote{\dagger}{The
ease of availability of different fonts is determined by the power of the
local {\tt texprint} program. A ``font manager'' to find obscure fonts is a
necessity if a wide range of fonts are needed \dots\ or desired.}
Fonts are available as members of a ``family'' or individually. 
Font families consist of various styles (normal) Roman, Boldface,
Italic, Slanted, and Mathematics symbols at the different point sizes. 

\shead{fontfamily}{Fonts in a Font Family}
A font family consists of a number of fonts and mathematical symbols that
have been chosen to satisfy the size requirements of sub and subsub scripts.
The family sizes available are 
{\twelvepoint |\twelvepoint|}, {\tenpoint |\tenpoint|}, 
{\ninepoint |\ninepoint|}, 
and 
{\eightpoint |\eightpoint|}.\footnote{\ddagger}{In \tex\ there are 72.26666 points to the inch (about 28.45 points to the
cm).}
To call a specific style of font in a font
family it is necessary to use the simple abbreviation of |\rm| for roman,
|\it| for text italic, |\sl| for slanted, |\bf| for bold, and |\tt| for a
fixed spacing typewriter like font. |\tt| is in the tenpoint typewriter font.
In addition, a font family has three {\it big} fonts called respectively,
|\bigfont|, |\biggfont|, and |\bigggfont|. These are all roman fonts and are
approximately 1.2, 1.44, and 1.78 times the basic family point size. There is
also a {\smallerfont smallerfont} that is  called through |\smallerfont|.
These fonts are all changed when the point size family is changed.

\shead{documentfonts}{Document Font}
\intex\ has the important idea of a {\it document} and along with it a 
{\it document font family}. 
This is the
default font for the document and the basis of the sizes of the fonts found
in the headers, footers, section heads, or footnotes among others. For
example the section head font is called |\sheadfont| and is only changed when
the document font family is changed and not by merely calling a point size
family. 
The change in the document font is made naming a particular {\tt <point size
font family>} {\bf as an argument} of a |\documentstyle| command. For
instance
|\documentstyle{\twelvepoint}| would create a document with the 
text twelve point Roman.
The fonts for
chapter and section heads, footnotes, headers, and footers would
also be appropriately modified. This
feature means that you could run |\eightpoint| text through several sections
without modifying the section head fonts. On the other hand, a change in
document font will change the entire document and can be done at any time or
even within a group. 

|\documentstyle{\tenpoint}| set the document font set of this book to tenpoint.

\font\fivermsca=cmr5 scaled 2074
\shead{specialfonts}{Additional Fonts} 
A new font, assuming that it exists, must first be defined and then called.
For example, |\font\fivermsca = cmr5 scaled 2074| {\bf defines} a five point
Roman font magnified by 2.074, or to about 10 points. The ``5 point'' is its
design size. To call the font, use |\fivermsca|, usually inside a group to
limit its effect. For instance ``{\fivermsca this was produced by}'' the
following ``|{\fivermsca this was produced by}|''. Note that a scaled five
point font is wider than a ten point form. The command is 

\beginblockmode
\mbr
\ext\@|\font\<fontname> [scaled <number>]|
\nbr
This defines a new font with the name |\<fontname>|. The scaling is optional.
The scale values are 4 digit integers, 1000 times greater than the nominal
size. Thus |\font\fvx = ambx5 scaled 1000| is identical to 
|\font\fvx = ambx5|. To actually use the font, the command |\<fontname>| is
called. 
\endblockmode
Some examples are as follows:
\font\tensc=cmcsc10
\font\twelvess=cmss12
\font\eighteenrm=cmr17
\font\tendu=cmdunh10
\font\tenu=cmu10
\font\sixrm=cmr6
\font\sixbf=cmbx6
\begingroup
\lineskip=3pt
\bshortcomlist
\ext\@|\fivei|&{\fivei Five point math italic}\cr
\ext\@|\fiverm|&{\fiverm Five point Roman}\cr
\ext\@|\fivebf|&{\fiverm Five point bold}\cr
|\font\sixrm=cmr6|&{\sixrm Six point Roman}\cr
|\font\sixbf=cmbx6|&{\sixbf Six point Bold}\cr
|\font\eighteenrm=cmr17|&{\eighteenrm Eighteen (Seventeen?) point Roman}\cr
|\font\tendu=cmdunh10|&{\tendu Ten point Dunhill}\cr
|\font\twelvess=cmss12|&{\twelvess Twelve point special San Serif}\cr
|\font\tenu=cmu10|&{\tenu Ten point unslanted italic}\cr
\eshortcomlist
\endgroup
Note that |\fivei| is a mathematics italic font. In this font the space has
been defined to have zero width and the character set is different than that
found in a typical Roman type font. 

\def\lb{\hfil\break}
\shead{fontcomlist}{Command Forms}
The four complete point size families are
\bshortcomlist
\ext\@|\twelvepoint|&{\twelverm Calls the twelvepoint font family}\cr
\ext\@|\tenpoint|&{\tenrm Calls the tenpoint font family}\cr
\ext\@|\ninepoint|&{\ninerm Calls the ninepoint font family}\cr
\ext\@|\eightpoint|&{\eightrm Calls the eightpoint font family}\cr
\eshortcomlist

\dssshead{Point Size Relative Fonts}
\begingroup
\veryraggedright
\bshortcomlist
\ext\@|\bf|&{\bf Bold in current font family  (tenpoint)}\cr
\ext\@|\it|&{\eightpoint \it Italic in current font family (eightpoint)}\cr
\ext\@|\rm|&{\eightrm Roman in current font family (eightpoint)}\cr
%\ext\@|\sl|&{\eightpoint \sl Slanted in current font family (eightpoint)}\cr
\ext\@|\tt|&{\eightpoint\tt Typewriter in current font family  (eightpoint)} \cr
\ext\@|\bigggfont|&{\eightpoint \bigggfont Biggg font in current  (eightpoint)}\cr
\ext\@|\biggfont|&{\eightpoint \biggfont Bigg font in current (eightpoint)}\cr
\ext\@|\bigfont|&{\eightpoint \bigfont Big font in current (eightpoint)}\cr
\ext\@|\cheadfont|&{\cheadfont chapterhead in current (tenpoint)}\cr
\ext\@|\documentstyle{<font family>}|& Sets document font family to {\tt <font family>}\cr
\ext\@|\footerfont|&{footerfont of current document font}\cr
\ext\@|\footnotefont|&{\footnotefont footnote font (family) of current document font} \cr
\ext\@|\headerfont|&{header font of current document font}\cr
\ext\@|\smallerfont|&{\smallerfont this is smaller relative to tenpoint}\cr
\ext\@|\sheadfont|&{\sheadfont Section head font of current document font}\cr
\ext\@|\ssheadfont|&{\ssheadfont Subsection head font of current document font}\cr
\ext\@|\sssheadfont|&{\sssheadfont Subsubsection head font of current document font}\cr
\ext\@|\dsssheadfont|&{\dsssheadfont Diminished section head font of current document font}\cr
\ext\@|\captionnumfont|&{\captionnumfont Caption label and number font of 
                                          current document font}\cr
\ext\@|\captiontitlefont|&{\captiontitlefont Caption title font of 
                                          current document font}\cr
\ext\@|\captionbodyfont|&{\captionbodyfont Caption body font of 
                                          current document font}\cr

\eshortcomlist
\endgroup

\shead{newfontfam}{New Font Families}
It is easy to produce a new font family that is compatible with \intex\---
just modify the |\tenpoint| definition that is in the {\tt
inrsfont.tex} file. 

\ejectpage
\done
